%% ---------------------------------------------------------------------
%% Chapter 1 -- Planning for Drone Swarms
%% The opening chapter of the book: the problem, the vocabulary, a preview
%% of the four-layer hybrid architecture, the map of the book, established
%% results versus open research, and how to read the book.
%% Code:    code/ch01_scenario.py (scenario generator with self-test)
%% Figures: figures/ch01/*.tex; scenario.tex is written by
%%          code/figures/gen_ch01_scenario.py (seed 1462).
%% ---------------------------------------------------------------------
\chapter{Planning for Drone Swarms}
\label{ch:ch01}

\begin{objectives}
\begin{itemize}
  \item Describe the problem that this book solves: a swarm of controlled drones
        flying planned, time-indexed routes through a shared space in which unknown,
        non-cooperative moving objects appear.
  \item Tell a \emph{planned conflict} between two drones of the swarm from an
        \emph{unexpected obstacle}, and say which of the two is resolved before
        take-off and which one during the flight.
  \item Use the vocabulary of the field correctly: workspace and configuration space,
        discrete and continuous planning, global and local planning, deliberative and
        reactive, centralized and decentralized, cooperative and non-cooperative,
        tracking and prediction, and the properties completeness, optimality, bounded
        suboptimality, anytime, incremental and real-time.
  \item Name the four layers of the hybrid architecture, say what each layer receives
        and delivers, and give the time scale on which each one works.
  \item Locate every algorithm of the training plan in this book, state its priority,
        and say which parts of the field are established and where research questions
        remain.
  \item Choose a reading path through the book and set up the Python stack that its
        code uses.
\end{itemize}
\end{objectives}

% =====================================================================
\section{The problem: many drones, one airspace, and a stranger}
\label{sec:ch01-problem}

Imagine that you are responsible for a fleet of eight quadrotors that inspect the
roof of a large factory every morning. Each drone has its own list of places to
photograph, and all of them share the same airspace. Before take-off you need a route
for every drone, and you need to be sure that no two routes bring two drones to the
same place at the same time. During the flight something unplanned happens: a drone
that is not yours crosses the site. It does not talk to your fleet, it does not follow
your plan, and you cannot know exactly where it is going. Your drones must notice it,
guess where it is heading, get out of its way, and then return to the plan---or, if
that is impossible, make a new plan and check it against the rest of the fleet.

This book is about exactly that problem, and this chapter states it carefully. Every
later chapter teaches one algorithm that solves one piece of it. The setting has four
ingredients.
\begin{itemize}
  \item \textbf{A swarm of controlled drones.}\index{swarm}\index{drone!controlled}
        There are $k$ drones that we command. We know their positions, we can give each
        of them a route, and they do what we tell them. In this book they are called
        the \emph{agents} of the swarm, or simply drones $A$, $B$, $C$, \dots
  \item \textbf{A shared space, represented on a grid or a graph.}\index{grid}\index{graph}
        The drones fly through one three-dimensional volume. For planning we
        discretise that volume into cells of a grid, or more generally into the
        vertices of a graph $G=(V,E)$ whose edges are the moves that a drone can make in
        one time step. \Cref{ch:ch02} explains the representations in detail; most
        figures of the book show a horizontal slice of the space, that is, a
        two-dimensional grid at the altitude of the swarm.
  \item \textbf{Planned, time-indexed routes.}\index{time-indexed path}
        A route is not just a curve. It says \emph{where} a drone is at \emph{each}
        time step. Only with time attached can we say whether two routes collide, and
        only then can a planner remove the collision before anybody takes off.
  \item \textbf{Unknown, non-cooperative moving objects.}\index{intruder}
        During the flight, objects that are not part of the swarm may appear: other
        drones, most often. We do not know their plans. We observe them with noisy
        sensors, so even their current state is uncertain, and their future motion is
        more uncertain still. The book calls such an object an \textbf{intruder}.
\end{itemize}

\begin{figure}[tb]
  \centering
  \inputfigure{ch01/scenario}
  \caption{The scenario of \cref{ex:ch01-scenario}: a $12\times 8$ grid with
    $14$ blocked cells, three controlled drones with their independently planned
    shortest paths, and one intruder entering from the top right with a velocity of
    $(-0.54,-0.38)$ cells per step. Two things go wrong. The paths of $A$ and $C$ both
    pass through cell $(3,3)$ at $t=4$ (red square): a \emph{planned conflict}, visible
    before take-off. The intruder's predicted motion (dotted, with growing uncertainty
    discs) passes within $0.38$ cells of drone $B$ at $t=6$ (dashed circle): an
    \emph{unexpected obstacle}, visible only during the flight.}
  \label{fig:ch01-scenario}
\end{figure}

\subsection{Two kinds of trouble}
\label{sec:ch01-two-kinds}

\Cref{fig:ch01-scenario} shows a small instance of the problem, and it shows the two
kinds of trouble that this book is about. The first kind is visible before anybody
takes off. Drone $A$ wants to fly from $(7,3)$ to $(0,0)$ and drone $C$ from $(6,4)$ to
$(3,1)$. Each of them, planning alone, chooses a shortest path, and both shortest paths
use cell $(3,3)$ at time step $4$. Nothing is unknown here: we know both routes, we
know both clocks, and we can see the collision coming while the drones are still on
the ground. The second kind of trouble is invisible before take-off. The intruder in
the top right corner was not in anybody's plan. Once it is observed, we can estimate
its velocity and extrapolate: if it keeps flying straight, it will pass within $0.38$
cells of drone $B$ at $t=6$. But it may also turn, slow down or stop, and the further
we look ahead, the less we know. The two kinds are defined next.

\begin{definition}[Planned conflict]
\label{def:ch01-planned-conflict}
A \textbf{planned conflict}\index{planned conflict|textbf} is a collision between two
controlled drones that follows from their planned routes alone: at some time step
both routes occupy the same vertex, or they traverse the same edge in opposite
directions. A planned conflict can be detected, and resolved, before the flight, by a
\emph{multi-agent planner} that knows all routes.
\end{definition}

\begin{definition}[Unexpected obstacle]
\label{def:ch01-unexpected-obstacle}
An \textbf{unexpected obstacle}\index{unexpected obstacle|textbf} is an object that
was not part of the plan and that is observed only during the flight. Its current
state is known only through noisy measurements and its future motion is uncertain. It
must be handled \emph{during} the flight, by whatever the drone can decide in the time
it has.
\end{definition}

\begin{figure}[tb]
  \centering
  \inputfigure{ch01/planned-vs-unexpected}
  \caption{The two kinds of trouble on a $5\times 5$ grid. (a)~Drones $A$ and $C$
    both want cell $(2,2)$ at $t=2$: a planned conflict. Everything about it is
    known before take-off, and letting $C$ wait one step removes it. (b)~An intruder
    whose plan we do not know: from its observed velocity we predict where it will be
    at $t=1,2,3$, with a disc of uncertainty that grows with the horizon, and the
    prediction crosses the route of $A$ near $t=3$. Nothing about it could have been
    planned for.}
  \label{fig:ch01-planned-vs-unexpected}
\end{figure}

The difference matters because the two kinds of trouble call for different tools.
A planned conflict is a problem of \emph{coordination} among agents that we fully
control; it is solved once, before the flight, by a planner that may take a few
seconds and that can guarantee a collision-free plan. An unexpected obstacle is a
problem of \emph{reaction} under uncertainty; it must be solved many times per second,
by each drone on its own, and no algorithm can guarantee safety against an object
whose motion is unknown. \Cref{fig:ch01-planned-vs-unexpected} contrasts the two
side by side.

\begin{keyidea}[The core idea of the book]
Use a \emph{global planner} for the planned routes, a \emph{local collision-avoidance
method} for the surprises, and a \emph{replanner} only when the local response is
insufficient.\index{hybrid architecture!core idea}
\end{keyidea}

This one sentence is the plan of the whole book. Chapters \ref{ch:ch03} to
\ref{ch:ch11} build the global planner, from single-agent graph search to multi-agent
path finding. Chapters \ref{ch:ch12} to \ref{ch:ch15} build the local layer. Chapters
\ref{ch:ch18} to \ref{ch:ch20} build the tracking and prediction that the local layer
needs in order to see the surprise coming. Chapters \ref{ch:ch05}, \ref{ch:ch16},
\ref{ch:ch17} and \ref{ch:ch21} to \ref{ch:ch23} provide the replanning, the
continuous-space planning and the control that hold everything together, and
\cref{ch:ch24,ch:ch25} assemble and evaluate the result.

\subsection{A scenario in numbers}
\label{sec:ch01-example}

Numbers make the distinction concrete. The scenario of \cref{fig:ch01-scenario} was
produced by the scenario generator \code{code/ch01\_scenario.py} of this chapter with
the fixed seed $1462$; every number quoted below is printed by that program.

\begin{example}[Three drones and one intruder]
\label{ex:ch01-scenario}
The grid has $12\times 8$ cells, $14$ of which are blocked. Cell $(x,y)$ covers the unit
square with lower-left corner $(x,y)$, its centre is $(x+0.5,\,y+0.5)$, and the grid is
4-connected: in one time step a drone moves to a horizontally or vertically adjacent
free cell, or it waits. A drone that has reached its goal stays there.
\begin{itemize}
  \item Drone $A$ flies from $(7,3)$ to $(0,0)$. Its shortest path has $10$ steps:
        $(7,3),(6,3),\dots,(1,3),(1,2),(1,1),(1,0),(0,0)$.
  \item Drone $B$ flies from $(1,5)$ to $(10,5)$ along row $5$ in $9$ steps.
  \item Drone $C$ flies from $(6,4)$ to $(3,1)$ in $6$ steps:
        $(6,4),(5,4),(4,4),(3,4),(3,3),(3,2),(3,1)$.
\end{itemize}
Planned independently, the three paths have a sum of costs of $10+9+6=25$ steps and a
makespan of $10$ steps (\cref{ch:ch02} defines both measures). They contain exactly one
planned conflict: a vertex conflict of $A$ and $C$ in cell $(3,3)$ at $t=4$. If either
drone waits one step before entering $(3,3)$, the conflict disappears and the sum of
costs rises to $26$.

The intruder is first observed at $(11.0,\,7.5)$ with velocity $(-0.54,\,-0.38)$ cells
per step, that is, a speed of $0.66$ cells per step. If it keeps that velocity, then
at $t=6$ it is at $(7.76,\,5.22)$ while drone $B$ occupies cell $(7,5)$ with centre
$(7.5,\,5.5)$: a distance of $0.38$ cells, far less than the safety distance that any
real swarm would demand. \Cref{tab:ch01-trace} lists the whole time line.
\end{example}

\begin{table}[tb]
  \centering\small
  \caption{Time line of \cref{ex:ch01-scenario}: the cell of each drone, the predicted
    position of the intruder, and its distance to the centre of each drone's cell (in
    cells). The planned conflict is the row in which $A$ and $C$ share a cell; the
    encounter is the row in which $d_B$ drops to $0.38$. Drones stay at their goals
    after arriving.}
  \label{tab:ch01-trace}
  \begin{tabular}{@{}rcccrrrr@{}}
  \toprule
  $t$ & $A$ & $B$ & $C$ & Intruder $(x,y)$ & $d_A$ & $d_B$ & $d_C$\\
  \midrule
  0 & (7,3) & (1,5) & (6,4) & (11.00, 7.50) & 5.32 & 9.71 & 5.41\\
  1 & (6,3) & (2,5) & (5,4) & (10.46, 7.12) & 5.37 & 8.12 & 5.61\\
  2 & (5,3) & (3,5) & (4,4) & (9.92, 6.74) & 5.48 & 6.54 & 5.86\\
  3 & (4,3) & (4,5) & (3,4) & (9.38, 6.36) & 5.66 & 4.96 & 6.17\\
  4 & \textbf{(3,3)} & (5,5) & \textbf{(3,3)} & (8.84, 5.98) & 5.89 & 3.37 & 5.89\\
  5 & (2,3) & (6,5) & (3,2) & (8.30, 5.60) & 6.17 & 1.80 & 5.71\\
  6 & (1,3) & (7,5) & (3,1) & (7.76, 5.22) & 6.49 & \textbf{0.38} & 5.66\\
  7 & (1,2) & (8,5) & (3,1) & (7.22, 4.84) & 6.18 & 1.44 & 5.00\\
  8 & (1,1) & (9,5) & (3,1) & (6.68, 4.46) & 5.97 & 3.01 & 4.34\\
  9 & (1,0) & (10,5) & (3,1) & (6.14, 4.08) & 5.86 & 4.59 & 3.69\\
  10 & (0,0) & (10,5) & (3,1) & (5.60, 3.70) & 6.02 & 5.22 & 3.04\\
  \bottomrule
  \end{tabular}
\end{table}

Look at what the table does and does not tell you. The conflict in row $t=4$ is a
fact: given the two routes, it will happen. The encounter in row $t=6$ is a
\emph{prediction}: it assumes that the intruder flies straight at constant velocity,
which is the simplest model we could adopt and, as \cref{ch:ch20} shows, not a bad one
over a short horizon. The dashed discs in \cref{fig:ch01-scenario} grow with $t$ to
remind you that the prediction gets worse the further ahead it looks. A good swarm
planner acts on the conflict now and on the encounter only when it is close enough in
time to be believed---which is the subject of the safety horizon in
\cref{sec:ch01-architecture}.

\begin{pitfall}[An intruder is not a wall]
A tempting shortcut is to treat an intruder as a static obstacle: block the cells it
currently occupies and replan around them. This fails twice. The intruder will have
moved by the time the new plan is executed, so the plan avoids a place where nothing
is, and it may steer the drone straight into where the intruder is heading. A moving
object must be avoided in \emph{space and time}, using a prediction of where it will
be, not a snapshot of where it was.\index{pitfall!static intruder}
\end{pitfall}

\subsection{Why not one method for everything?}
\label{sec:ch01-why-hybrid}

You may ask why the book does not simply pick the most powerful planner and run it
whenever anything changes. Two arguments answer this, and they recur throughout the
book. First, \emph{time}. A multi-agent planner that guarantees a collision-free plan
for the whole swarm needs seconds, sometimes much longer, on instances of realistic
size. A drone flying at a few metres per second cannot wait that long when another
drone is two metres away. Second, \emph{guarantees}. The global plan has properties that
we paid for: it is collision-free among the swarm, it is optimal or provably close to
optimal, and it respects formation and communication constraints. Throwing it away at
every surprise means paying again, and possibly losing those properties in the
meantime. The opposite extreme is equally poor: a swarm that only reacts locally
never plans, and \cref{exr:ch01-deadlock} asks you to show that such a swarm can lock
up in a corridor with nobody willing to go first. The hybrid architecture previewed in
\cref{sec:ch01-architecture} keeps the global plan as the default, reacts locally when
it must, and replans only when the local reaction cannot bring the drone back.

% =====================================================================
\section{Vocabulary: the words you need before anything else}
\label{sec:ch01-vocabulary}

Planning has a vocabulary of contrasts. Most of them are pairs of words that describe
two ends of a spectrum, and every algorithm in this book sits somewhere between them.
\Cref{tab:ch01-contrasts} collects the pairs; the paragraphs below explain each one
once, and the rest of the book uses them without further comment.

\begin{table}[tb]
  \centering\small
  \caption{The contrasts that organise the book. Most algorithms sit at one end of
    each pair; the hybrid architecture of \cref{ch:ch24} deliberately combines both
    ends.}
  \label{tab:ch01-contrasts}
  \begin{tabularx}{\textwidth}{@{}lYY@{}}
  \toprule
  Pair & One end & The other end\\
  \midrule
  Workspace / configuration space & the physical volume the drone flies in & the set of all configurations of the drone; obstacles are inflated by its size\\
  Discrete / continuous & finite vertices and time steps; graph search & real-valued positions and time; sampling, optimisation\\
  Global / local & whole route from start to goal, all obstacles known & next few seconds, only what the sensors see\\
  Planning / control & decides \emph{where} to go and \emph{when} & makes the vehicle actually follow the decision\\
  Deliberative / reactive & thinks ahead, uses a model, takes time & maps the current situation directly to an action\\
  Centralized / decentralized & one computer knows and decides everything & every drone decides from what it knows and hears\\
  Cooperative / non-cooperative & the other agent shares the avoidance rules & the other agent does what it likes\\
  Tracking / prediction & estimating the state \emph{now} from noisy data & extrapolating the state into the future\\
  \bottomrule
  \end{tabularx}
\end{table}

\paragraph{Workspace and configuration space.}
The \textbf{workspace}\index{workspace} is the physical space in which the drone
moves: a volume of $\R^3$, or a horizontal slice of it. The \textbf{configuration
space}\index{configuration space} $\Cspace$ is the set of all configurations of the
drone; for a drone that we model as a point with a safety radius, a configuration is
its position, and the obstacle region $\Cobs$ is the workspace obstacle inflated by
that radius. Planning takes place in $\Cspace$, so that a path of the point robot in
$\Cfree=\Cspace\setminus\Cobs$ is a collision-free motion of the real drone. That is
all you need for now; \cref{ch:ch02} gives the definitions, the inflation operation
(the Minkowski sum) and the usual approximations \cite{lavalle2006planning}.

\paragraph{Discrete and continuous planning.}
\textbf{Discrete planning}\index{planning!discrete} works on a graph: a finite set of
vertices, a finite set of moves, and integer time steps. Dijkstra's algorithm, \astar,
\dstarlite and all the multi-agent planners of the book are discrete.
\textbf{Continuous planning}\index{planning!continuous} works with real-valued
positions, velocities and times; the sampling-based planners \rrt and \rrtstar and the
optimisation-based methods MPC and MILP live there, and so do the velocity-space
methods of local avoidance. The two are connected by discretisation: a grid is a
sampling of the continuous workspace, and a discrete plan becomes a continuous
trajectory when the drone flies between the cell centres.

\paragraph{Time-indexed paths.}
The single most important object in this book is a path with time attached.

\begin{definition}[Time-indexed path]
\label{def:ch01-time-indexed-path}
A \textbf{time-indexed path}\index{time-indexed path|textbf} of an agent on a graph
$G=(V,E)$ is a sequence $\pi=(v_0,v_1,\dots,v_T)$ of vertices such that $v_0$ is the
agent's start, $v_T$ is its goal, and for every $t<T$ either $v_{t+1}=v_t$ (a
\emph{wait} action) or $\set{v_t,v_{t+1}}\in E$ (a \emph{move} action). The agent
occupies $v_t$ during time step $t$ and stays at $v_T$ for all $t\ge T$. A
\textbf{plan}\index{plan} for a swarm is one time-indexed path per agent.
\end{definition}

Without the index $t$ two paths can cross without a collision, because the drones pass
the crossing at different times; with the index, a collision is a precise statement:
$v_t^A=v_t^C$ for some $t$. \Cref{ch:ch02} turns this into the space-time (or
time-expanded) graph and \cref{ch:ch07} lists every kind of conflict that time-indexed
paths can have \cite{stern2019mapf}.

\paragraph{Global and local planning.}
A \textbf{global planner}\index{planning!global} computes a complete route from the
start to the goal, using a map of all known obstacles; it is allowed to take time and
it may be optimal. A \textbf{local planner}\index{planning!local} looks only a few
seconds ahead, uses only what the sensors currently see, and outputs the next
velocity or the next short manoeuvre. Global planners know the goal but not the
surprise; local planners know the surprise but not the goal. The book's local methods
are velocity obstacles, \orca, the dynamic window approach and potential fields
(\cref{ch:ch12,ch:ch13,ch:ch14,ch:ch15}); its global methods are everything in
\cref{ch:ch03,ch:ch04,ch:ch05,ch:ch06,ch:ch07,ch:ch08,ch:ch09,ch:ch10,ch:ch11,ch:ch16,ch:ch17}.

\paragraph{Planning and control.}
\textbf{Planning}\index{planning} decides \emph{where} the drone should be and
\emph{when}; \textbf{control}\index{control} makes the physical vehicle follow that
decision by computing motor commands from the error between the desired and the
measured state. The book plans; it treats the flight controller as a black box that
receives a velocity or a waypoint and tracks it. The exception is model predictive
control (\cref{ch:ch21}), which plans and controls at once by re-solving a small
optimisation problem at every control cycle---which is why it is such a natural place
to enforce constraints.

\paragraph{Deliberative and reactive.}
A \textbf{deliberative}\index{deliberative} method thinks ahead: it uses a model of the
world to search or optimise over possible futures before it acts. A
\textbf{reactive}\index{reactive} method maps the current situation to an action with
little or no search, which makes it fast and robust to modelling errors but
short-sighted. The two words are the behavioural counterparts of global and local:
global planners are deliberative, local avoidance is reactive, and the hybrid
architecture layers the one on the other \cite{russell2020aima}.

\paragraph{Centralized and decentralized.}
A \textbf{centralized}\index{centralized} system has one computer that knows the state
of every drone and decides for all of them; a \textbf{decentralized}\index{decentralized}
system lets every drone decide from its own observations and from what it hears from
its neighbours over a communication link of limited range. The global planner of the
book is centralized: conflict-based search needs all routes in one place. Local
avoidance is decentralized: each drone runs \orca on what it sees. Consensus and
formation control (\cref{ch:ch23}) are decentralized by design, and the communication
radius that limits them is one of the constraints that the research directions of
\cref{sec:ch01-research} are about.

\paragraph{Cooperative and non-cooperative agents.}
Two agents are \textbf{cooperative}\index{cooperative agent} if both follow the same
avoidance rules, so that each may count on the other to do its share; reciprocal
velocity obstacles and \orca are built on exactly this assumption
\cite{vandenberg2011orca}. A \textbf{non-cooperative}\index{non-cooperative agent}
agent does whatever it does; the intruder of \cref{fig:ch01-scenario} is one, and a
drone that meets it must take the full responsibility for avoiding it. The drones of
the swarm are cooperative with one another and non-cooperative with the intruder,
and \cref{ch:ch13} shows that the same algorithm handles both cases with a change of a
single factor.

\paragraph{Tracking and prediction.}
\textbf{Tracking}\index{tracking} estimates the \emph{current} state of a moving
object---position and velocity---from a sequence of noisy measurements; the Kalman
filter of \cref{ch:ch18} is the classical tool \cite{thrun2005probabilistic}.
\textbf{Prediction}\index{prediction} extrapolates the estimated state into the
\emph{future}, over a horizon of a few seconds, and reports the growing uncertainty
along the way (\cref{ch:ch20}). The two are often confused because the Kalman filter
contains a prediction step; the distinction is that tracking ends at the present
moment and prediction begins there.

\subsection{Properties of algorithms}
\label{sec:ch01-properties}

When the book says that an algorithm is complete, or optimal, or anytime, it means
something precise. The precise statements come with the algorithms; here are the
informal meanings, so that you can read the map of the book in
\cref{sec:ch01-map} with them in mind. \Cref{tab:ch01-properties} summarises them.

\begin{table}[tb]
  \centering\small
  \caption{Properties of planning algorithms, informally. The formal versions appear
    as theorems in the chapters listed in the last column.}
  \label{tab:ch01-properties}
  \begin{tabularx}{\textwidth}{@{}lYl@{}}
  \toprule
  Property & Informal meaning & Defined in\\
  \midrule
  Complete & if a solution exists, the algorithm finds one (and if none exists, it says so) & \cref{ch:ch03,ch:ch09}\\
  Optimal & the solution it returns has the least possible cost & \cref{ch:ch04,ch:ch09}\\
  Bounded suboptimal & the cost is at most $w$ times the optimal cost, for a factor $w\ge 1$ chosen by the user & \cref{ch:ch04,ch:ch10}\\
  Anytime & returns a first solution quickly and improves it while time remains & \cref{ch:ch06}\\
  Incremental & reuses the previous search when the problem changes slightly, instead of starting over & \cref{ch:ch05}\\
  Real-time & is guaranteed to return an answer within a fixed, short time budget & \cref{ch:ch13,ch:ch14}\\
  \bottomrule
  \end{tabularx}
\end{table}

\begin{description}[style=nextline,leftmargin=1.5em]
  \item[Completeness.]\index{completeness} An algorithm is \textbf{complete} if it
    finds a solution whenever one exists, and reports failure when none exists. On a
    finite graph, Dijkstra's algorithm and \astar are complete. Prioritized planning is
    not: it may fail on an instance that has a solution (\cref{ch:ch08}). Sampling-based
    planners are \emph{probabilistically} complete: the probability that they find a
    solution, if one exists, tends to one as the number of samples grows
    (\cref{ch:ch16}).
  \item[Optimality.]\index{optimality} An algorithm is \textbf{optimal} if the solution
    it returns has the least cost among all solutions. \astar with an admissible
    heuristic is optimal (\cref{ch:ch04}); conflict-based search is optimal for the
    sum of costs of a multi-agent plan (\cref{ch:ch09}); \rrtstar is
    \emph{asymptotically} optimal, that is, optimal in the limit of infinitely many
    samples (\cref{ch:ch17}).
  \item[Bounded suboptimality.]\index{bounded suboptimality} An algorithm is
    \textbf{bounded suboptimal} with factor $w\ge 1$ if the cost of its solution is at
    most $w$ times the optimal cost. It trades a known, chosen loss of quality for
    speed. Weighted \astar and \ecbs are the book's two examples
    (\cref{ch:ch04,ch:ch10}); with $w=1$ they become optimal, with $w=1.5$ they are
    often orders of magnitude faster.
  \item[Anytime.]\index{anytime algorithm} An \textbf{anytime} algorithm returns a
    first, possibly poor, solution quickly and then keeps improving it as long as it
    is allowed to run; interrupt it whenever you like and take the best solution so
    far. \arastar does this by lowering the factor $w$ step by step (\cref{ch:ch06}).
  \item[Incremental search (replanning).]\index{incremental search}\index{replanning}
    An \textbf{incremental} search reuses the work of the previous search when the
    graph changes a little---an obstacle appears, an edge cost changes---instead of
    starting from scratch. \lpastar and \dstarlite are the book's incremental
    planners (\cref{ch:ch05}); they are the replanning layer of the hybrid
    architecture.
  \item[Real-time.]\index{real-time} A \textbf{real-time} algorithm is guaranteed to
    return an answer within a fixed, short time budget, whatever the situation, because
    the amount of work it does per call is bounded. \orca solves a small linear
    program per drone per cycle; the dynamic window approach evaluates a fixed number of
    candidate velocities. Neither searches for a route, and that is precisely why they
    can run at tens of hertz.
\end{description}

These properties conflict with one another, and choosing among them is the designer's
job. An optimal multi-agent planner is not real-time; a real-time local method is not
complete; an incremental planner can only repair what a complete planner has built.
The four layers of the next section are chosen so that each property is provided by
the layer that can afford it.
